\section{Closest native formalisms and the redundancy test}

The most important literature for ICP is not the literature that sounds most
similar. It is the literature capable of making the proposal redundant. Five
clusters cover much of the generic future-constraint language: control and
viability, the adjacent possible, empowerment and causal information,
constructor and predictive-state theory, and feedback information
thermodynamics.

\subsection{Reachability and viability already describe constrained futures}

Hamilton--Jacobi reachability represents controlled future sets directly.
Mitchell, Bayen, and Tomlin characterize a finite-horizon backward-reachable set
for a nonlinear differential game as a sublevel set of a viscosity solution to
a Hamilton--Jacobi--Isaacs equation \citep{MitchellEtAl2005Reachable}. Viability
theory identifies states from which at least one admissible evolution can remain
inside a declared constraint set, with convergent approximation schemes under
stated conditions \citep{SaintPierre1994Viability}. These are not loose
analogies to Eq.~\eqref{eq:future-family}; they are mature mathematical theories
of accessible and survivable futures under controls and disturbances.

The distinction available to ICP lies upstream of the computed set. A
reachability calculation is given the state, law, targets, controls,
disturbances, boundary, and horizon. It need not ask how a particular
constraint role was formed, where it is physically realized, or whether it can
survive replacement of its first witness. If ICP stops at ``the state changes
what is reachable,'' it is standard control language. If it can identify a
minimal role that is formed and transported across realizations, it may add a
new state-equivalence or mechanism question to that language.

\subsection{The adjacent possible is the closest support-creation prior art}

Kauffman defines the chemical adjacent possible as molecular species absent
from the actual but synthesizable from it in one reaction step
\citep[p.~142]{Kauffman2000Investigations}. In the notation of
Eq.~\eqref{eq:adjacent-frontier}, this is a novelty-filtered, type-projected
one-step reachable set. That overlap is unusually strong: the adjacent possible
is native prior art for the part of ICP concerned not only with excluding
futures but with making new outcome types accessible.
Kauffman's later BioSystems formulation makes the stronger claim that enabling
constraints create unprestatable adjacent possibilities in biological
evolution \citep{Kauffman2014Prolegomenon}; a recent extension contrasts fixed
law-governed possibility spaces with an indefinite ``domain of no law''
\citep{PatraKauffman2025NoLaw}. These are important ontology-level competitors,
not empirically established consequences of the one-step definition.

The closest geometric competitor is still more direct. Fari\~na and Vallejo
define a ``biological cosmological constant'' $\Lambda_B$ proportional to the
squared fractional expansion rate of Fisher-information volume, interpret the
origin of life as a geometric phase transition, and apply the construction to
staged Phanerozoic diversification \citep{FarinaVallejo2026LambdaB}. This is
explicit prior art for geometric expansion into the adjacent possible, not
evidence for ICP. Its eco-evolutionary manifold and nonnegative squared rate do
not identify a replaceable carrier, preserve enabled/lost/retained identities,
or distinguish signed local pruning from enabling. ICP therefore cannot claim
the first geometric adjacent-possible measure or a generic phase-transition
signature; its narrower opportunity is an intervention-defined signed contrast
and a prospectively predicted carrier-linked crossover.

Later work turns the idea into several non-equivalent models. The urn model
with triggering adds $\nu+1$ previously unseen colors when a novel color is
drawn, while reinforcement $\rho$ increases reuse; it yields correlated
novelty together with Heaps- and Zipf-like scaling in the stated regimes and
was compared with textual, music, tagging, and Wikipedia sequences
\citep{TriaEtAl2014Novelties}. Edge-reinforced walks instead place discovery on
a declared concept network \citep{IacopiniEtAl2018Innovation}; multi-agent urns
model evolving social contacts \citep{UbaldiEtAl2021SocialNetworks}; and
higher-order novelty counts first occurrences of relations or combinations,
which can distinguish sequences with similar first-order discovery rates
\citep{DiBonaEtAl2025HigherOrder}. Dynamic scaffolding can also change which
possibilities are developmentally available, although the published treatment
is a modeling proposal rather than a direct test
\citep{Griesemer2024Scaffolded}. Numerical work in a Kauffman chemical model
further reports that autocatalytic closure can causally change complexity
dynamics under declared interventions \citep{Faggian2025Closure}; causal
organization-to-novelty effects are therefore prior art as well as conceptual
expectations.

The combinatorial TAP family makes the expansion explicit. One studied form is
\begin{equation}
 M_{t+1}=(1-\mu)M_t+
 \sum_{i=2}^{M_t}\alpha_i\binom{M_t}{i},
 \label{eq:tap-equation}
\end{equation}
where subsets of $M_t$ existing types generate candidates and $\mu$ represents
loss. Stochastic and deterministic variants can produce a long plateau followed
by rapid or finite-time explosive growth
\citep{SteelEtAl2020Combinatorial,CortesEtAl2025TAP}. This behavior is not a
universal observational law. Explicit-type models show that duplicate
rediscovery and finite worlds matter \citep{VillaniSerra2023BinaryString}, and a
1908--2016 Swedish product study found resource-constrained innovation rates,
Heaps-like type discovery, and useful but incomplete predictive information in
product-space topology rather than unrestricted combinatorial blow-up
\citep{Taalbi2026AdjacentPossible}.

This literature changes the novelty verdict. ICP cannot claim invention of
history-dependent accessibility, triggered support growth, novelty clustering,
scaffolded opening of opportunities, or a generic plateau-to-explosion
transition. It also cannot claim the first Fisher-geometric or cosmological
formalization of adjacent-possible expansion. Nor do Heaps or Zipf exponents
identify the mechanism. The possible
remaining addition is causal and physical: predeclare the type map, frontier,
carrier, horizon, and resources; alter or replace that carrier; and predict the
enabled, lost, and retained frontier on held-out interventions. Formal AP models
usually prestate an alphabet, graph, urn, or combinatorial generator. The
stronger claim that biological phase space is intrinsically unprestatable
\citep{LongoEtAl2012Enablement,Kauffman2014Prolegomenon,
PatraKauffman2025NoLaw} cannot simultaneously serve as a directly enumerated
experimental variable without an operational surrogate.

\subsection{Empowerment and future-option measures}

Empowerment makes the overlap still sharper. Its original finite-step form is
the channel capacity from present action sequences to later sensed states---a
scalar measure of potential perceivable control \citep{KlyubinEtAl2005Empowerment}.
Process empowerment replaces fixed open-loop scripts with feedback-conditioned
probing policies and can expose controllability hidden by noise or traps
\citep{TiomkinEtAl2025ProcessEmpowerment}. Causal entropic forcing instead
builds a finite-horizon path-entropy objective into an effective force and
produces adaptive-looking behavior in simple simulations
\citep{WissnerGrossFreer2013CausalEntropy}.

These works defeat any claim that ICP discovered the idea of valuing or
counting future options. They also clarify what ICP does \emph{not} assert. The
principle does not presently say that systems maximize reachable-set size,
empowerment, entropy, or optionality. A constraint may reduce, redirect, or
reorganize futures; persistent structure can increase viable options at one
scale while excluding paths at another. Scalar option values also collapse
outcome identity: equal capacities need not imply equal future-law families.

\subsection{Causal information and the closest semantic rival}

Causal-information measures require post-intervention distributions or a
defended causal model. Ay and Polani formalize information flow using causal
rather than merely observational independence \citep{AyPolani2008InformationFlows}.
Counterexamples show why transfer entropy and related dyadic quantities can
both overestimate flow and underestimate influence
\citep{JamesEtAl2016InformationFlows}. Predictive association, including a
Zipf law or high mutual information, is therefore not a carrier test.

Kolchinsky and Wolpert provide the closest mature formal rival to the phrase
``information that makes a difference.'' They define semantic information as
system--environment correlation whose controlled scrambling lowers a declared
future viability \citep{KolchinskyWolpert2018Semantic}. The intervention,
system boundary, ensemble, horizon, syntactic measure, and viability function
are constitutive. On a survival or self-maintenance task, an ICP analysis should
specialize or compare with this framework, not silently rename it.

ICP is broader only in a limited sense: its target may be any registered
future-response role, not necessarily viability, and it emphasizes formation
and succession across original-witness replacement. Breadth is not novelty by
itself. Unless carrier succession produces a new invariant or empirical
prediction, semantic information and ordinary causal-effect analysis may
already exhaust the phenomenon.

\subsection{Constructor theory and predictive-state theory}

Constructor theory supplies a law-level language of possible and impossible
repeatable tasks and characterizes information variables through physical task
repertoires such as copying, permutation, and distinguishability
\citep{Deutsch2013Constructor,DeutschMarletto2015Information}. This is the
cleanest neighbor for hard physical possibility. It does not, in its basic task
definition, supply a history-conditioned stochastic process or require that a
present constructor was formed through carrier succession. ICP may eventually
add a temporal formation and inheritance layer; for now the relationship is
structural, not derivational.

Computational mechanics begins from histories rather than tasks. Causal states
are equivalence classes of pasts with the same conditional distribution over
futures and yield a minimal optimal predictive representation under the stated
process assumptions \citep{ShaliziCrutchfield2001ComputationalMechanics}.
Input--output epsilon-transducers extend the construction to conditional output
laws under future input sequences \citep{BarnettCrutchfield2015Transducer}.
This is almost exactly the $H\rightarrow M\rightarrow P(F\mid M)$ part of
ICP.

The difference is physical, not algebraic. An inferred causal state need not be
localized, manipulable, or Pearl-causal; conditioning on an input sequence is
not automatically a surgical intervention. ICP's carrier criterion asks
whether a physically embodied realization is necessary and sufficient for the
relevant family and whether its replacement preserves the role. If that
physical test adds nothing to the predictive equivalence class, computational
mechanics is the more mature account.

\subsection{Feedback information thermodynamics}

Feedback thermodynamics gives direct evidence that a retained measurement can
change a later physical transformation. Sagawa and Ueda's generalized
Jarzynski relation includes information and feedback efficacy in work bounds
\citep{SagawaUeda2010Feedback}; Toyabe and colleagues tested information-to-free-
energy conversion with position-conditioned feedback on a Brownian particle
\citep{ToyabeEtAl2010InformationEnergy}. An autonomous bit-tape demon makes the
memory carrier and its thermodynamic accounting explicit
\citep{MandalJarzynski2012Demon}. Current work also shows that finite memory and
a limited protocol repertoire alter the optimal feedback value
\citep{HartleEtAl2026LimitedResources}.

A 2026 experimental preprint provides an especially relevant boundary case.
Time-delayed measurements of a colloid in a viscoelastic bath reveal memory not
captured by the visible instantaneous coordinate, and feedback can extract work
from hidden environmental degrees of freedom \citep{MurugaEtAl2026HiddenDegrees}.
This is strong evidence that information can persist outside the original or
currently observed bearer while remaining physically realized elsewhere. It
also supports the state-augmentation null: the enlarged system is Markovian and
the hidden modes carry ordinary excess free energy. As a preprint, it remains
frontier evidence.

Information thermodynamics therefore supports causal utility and closed
physical accounting. It does not support the inference that energy emerges from
information, because energy, Hamiltonians, baths, memories, and protocols enter
the theory as supplied structure. IOTA must derive mass- and energy-like
invariants at a later rung rather than cite feedback conversion as that
derivation.

\subsection{Result of the redundancy audit}

The generic ICP proposition is largely preempted. Mature theories already
establish that present physical states, controls, memories, and correlations
differentiate future possibilities and performance. The viable novelty target
is narrower:

\begin{quote}
An interaction forms a physically embodied constraint role; the original
witness ceases to instantiate it; a successor realization preserves the
registered intervention-response role; and perturbing that successor changes
the future family under matched controls.
\end{quote}

That target aligns with the strongest IOTA evidence and remains open as a
general principle. It is also falsifiable by complete-state reduction,
failure of role equivalence, or representation dependence.
